% !Mode:: "TeX:UTF-8"

\chapter{端到端量子接口仿真的张量网络方法}[Tensor-Network Methods for End-to-End Quantum Interface Simulation]

端到端量子接口链路同时包含连续时间量子光场、离散内部能级、有限带宽记忆以及条件测量。若直接在连续场与全密度矩阵层面处理，状态空间会随时间分辨率和链路长度迅速膨胀，数值规模很快失去控制。所研究方案之所以能够抽象为矩阵乘积态（matrix product state, MPS）与时间演化块消元（time-evolving block decimation, TEBD）框架，关键在于三点：第一，输入输出理论允许把连续时间输出场离散为一串按时间顺序排列的时间 bin \cite{GardinerCollett1985,Pichler2016TimeDelaysMPS}；第二，原子、时间 bin 与滤波记忆模之间的相互作用可以组织为一维链上的局域门序列 \cite{Pichler2016TimeDelaysMPS,Wang2017MPSCascadedQuantumSystems,ArranzRegidor2021WaveguideQED}；第三，中间站线性光学网络和探测器响应能够通过 Heisenberg 对偶映射推回为输入端的等效 effect，从而把“态演化”和“测量统计”分到两条互相兼容的计算通道中 \cite{Kraus1983StatesEffectsOperations,NielsenChuang2010QCQI}。

\section{连续时间链路的离散表示}[Discretizing the Continuous-Time Link]

量子接口实验中的光场天然是连续时间对象。原子在时刻 $t$ 向外场辐射，探测器也在一个有限时间窗内记录点击事件。若直接把场写成连续模，理论上当然最完整，但数值上既不利于构造局域张量，也不利于和实验中的离散时间记录对应。输入输出理论提供了一种更适合仿真的表述：在足够小的时间步长 $\Delta t$ 下，把连续场划分为一串彼此独立的时间 bin，每个 bin 只对应一个有限时间窗内的场自由度 \cite{GardinerCollett1985,Pichler2016TimeDelaysMPS}。

\subsection{时间 bin 的定义}[Definition of Time Bins]

设输出场算符为 $b_{\mu}(t)$，其中标记 $\mu$ 用来区分偏振或频率等模式自由度。连续场满足等时对易关系
\begin{equation}
[b_{\mu}(t),b_{\nu}^{\dagger}(t')]=\delta_{\mu\nu}\,\delta(t-t').
\label{eq:ch2-continuum-comm}
\end{equation}
把时间轴划分为 $t_k=k\Delta t$ 之后，可定义第 $k$ 个 bin 的离散化场算符
\begin{equation}
\Delta B_{k,\mu}=\frac{1}{\sqrt{\Delta t}}\int_{t_k}^{t_{k+1}} b_{\mu}(t)\,\mathrm d t,
\label{eq:ch2-bin-operator}
\end{equation}
于是有
\begin{equation}
[\Delta B_{k,\mu},\Delta B_{l,\nu}^{\dagger}]=\delta_{kl}\,\delta_{\mu\nu}.
\label{eq:ch2-bin-comm}
\end{equation}
式 \eqref{eq:ch2-bin-comm} 的意义很直接：在离散后的模型里，不同时间窗的场模已经可以当作一串独立站点。经过离散化后，连续时间光场可写成一维链上的一列局域自由度，原先的“无限模环境”描述被替换为顺序排列的时间站点 \cite{Pichler2016TimeDelaysMPS,ArranzRegidor2021WaveguideQED}。

单光子波包在离散化后也有清晰表达 \cite{Pichler2016TimeDelaysMPS}。若连续时间包络为 $\xi_{\mu}(t)$，则单光子态可写为
\begin{equation}
|1_{\xi}\rangle=\sum_{k,\mu}\sqrt{\Delta t}\,\xi_{\mu}(t_k)\,\Delta B_{k,\mu}^{\dagger}|\mathrm{vac}\rangle,
\label{eq:ch2-wavepacket-discrete}
\end{equation}
其归一化条件变为
\begin{equation}
\sum_{k,\mu}\Delta t\,|\xi_{\mu}(t_k)|^2=1.
\label{eq:ch2-wavepacket-norm}
\end{equation}
从实验角度看，式 \eqref{eq:ch2-wavepacket-discrete} 有两个直接好处。其一，波包包络、到达时间抖动和探测时间窗都可以在同一网格上比较。其二，时间分辨率由 $\Delta t$ 明确给出，离散误差也随之转化为可扫描、可收敛检查的数值参数 \cite{HongOuMandel1987,vanleent2022entangling}。

\begin{figure}[htbp]
  \centering
  \includegraphics[width=0.88\textwidth]{PichlerZoller2015_Fig1_timebin.png}
  \caption{连续时间光场离散为时间 bin 链的思想示意。系统在每一个时间步只与当前 bin 发生局域相互作用；当链路存在延迟或反馈时，先前发出的 bin 会在后续时刻再次参与动力学。图源引自 Pichler 和 Zoller 的时间延迟光子线路工作 \cite{Pichler2016TimeDelaysMPS}。}
  \label{fig:ch2-timebin}
\end{figure}

\subsection{一维链映射}[One-Dimensional Chain Mapping]

时间离散化之后，还需要把“真实链路中的许多物理对象”组织成适合张量网络的站点顺序。所研究链路的核心对象包括两臂发射体、传播中的时间 bin，以及 QFC 之后有限带宽滤波器引入的记忆模。将它们写成直积空间后，可得到
\begin{equation}
\mathcal H_{\mathrm{chain}}
=
\mathcal H_{\mathrm{em}}^{(A)}\otimes
\Bigl(\bigotimes_{k=1}^{N_t}\mathcal H_{k}^{(A)}\Bigr)\otimes
\mathcal H_{\mathrm{mem}}^{(A)}\otimes
\mathcal H_{\mathrm{em}}^{(B)}\otimes
\Bigl(\bigotimes_{k=1}^{N_t}\mathcal H_{k}^{(B)}\Bigr)\otimes
\mathcal H_{\mathrm{mem}}^{(B)}.
\label{eq:ch2-chain-space}
\end{equation}
在后续实现里，发射体站点采用 12 维局域空间，时间 bin 采用 5 维局域空间，滤波记忆模采用 3 维局域空间。这一局域维度选择有明确的物理对应关系：12 维发射体用于同时容纳原子内部态和腔偏振自由度，5 维 bin 用于表示 $|\varnothing\rangle$、$|H_{780}\rangle$、$|V_{780}\rangle$、$|H_{1517}\rangle$ 与 $|V_{1517}\rangle$，3 维记忆模用于表示真空、单激发以及有限带宽滤波保留下来的主要相干记忆。局域维度一旦确定，门算符的维度和张量收缩路径也随之固定下来。

这一映射最重要的物理依据是因果顺序。对于发射阶段，系统在时间步 $k$ 只需要和第 $k$ 个 bin 相互作用。对于 QFC 与滤波阶段，当前 bin 只与本地记忆模相互作用。对于中间站之前的光纤和线性光学网络，传播效应则可以在测量端通过 effect 推回处理 \cite{Pichler2016TimeDelaysMPS,Wang2017MPSCascadedQuantumSystems,ArranzRegidor2021WaveguideQED}。这种“每一步只接触有限邻域”的结构，使链路能够写成局域门序列，不必构造全局演化算符。

在离散时间步长足够小的情况下，发射体与当前 bin 的一次短时相互作用可写为
\begin{equation}
U_k^{\mathrm{(em)}}\approx
\exp\!\left[
-iH_{\mathrm{sys}}\Delta t
+
\sum_{\mu}
\sqrt{\Delta t}
\bigl(
L_{\mu}\,\Delta B_{k,\mu}^{\dagger}
-
L_{\mu}^{\dagger}\,\Delta B_{k,\mu}
\bigr)
\right],
\label{eq:ch2-emission-gate}
\end{equation}
其中 $H_{\mathrm{sys}}$ 描述发射体内部相干动力学，$L_{\mu}$ 描述向不同场模的辐射耦合。式 \eqref{eq:ch2-emission-gate} 只作用于“发射体 $+$ 当前 bin”两个相邻站点，因而天然适合 TEBD 的两站点更新。QFC 门与滤波记忆门具有同样的结构：虽然具体矩阵元素不同，但它们同样是局域站点上的幺正门或 Kraus 门。这正是第三章能够把连续时间接口链路压缩为“一维链 $+$ 局域门”的根本原因。

\section{矩阵乘积态表示}[Matrix-Product-State Representation]

把连续场离散成时间 bin 链之后，接下来的重点转为状态的存储和更新。若把整条链路看成由 $N$ 个局域站点组成、每个站点维度约为 $d$，则纯态可以形式上写成
\begin{equation}
|\psi\rangle=
\sum_{s_1,\ldots,s_N=1}^{d}
C_{s_1\cdots s_N}
|s_1\rangle\otimes\cdots\otimes|s_N\rangle,
\label{eq:ch2-full-state}
\end{equation}
其中 $C_{s_1\cdots s_N}$ 是 $N$ 阶复张量。直接存储式 \eqref{eq:ch2-full-state} 需要的参数量随 $N$ 指数增长，因此只有当态的相关结构非常特殊时，才可能存在低秩压缩。MPS 正是这种压缩在一维链上的标准表达，它的成立基础既来自 Schmidt 分解，也来自一维局域相互作用体系常见的有限纠缠结构 \cite{White1992DMRG,White1993DMRG,Vidal2003SlightlyEntangled,Schollwoeck2011MPS,Orus2014TensorNetworks}。

\subsection{Schmidt 分解和截断}[Schmidt Decomposition and Truncation]

对任意切分 $[1,\ldots,n]\,|\,[n+1,\ldots,N]$，把系数张量重排为矩阵 $M^{(n)}$ 并做奇异值分解（singular value decomposition, SVD），可得 Schmidt 展开
\begin{equation}
|\psi\rangle=
\sum_{\alpha=1}^{r_n}
\lambda_{\alpha}^{(n)}
|\alpha\rangle_{[1:n]}
\otimes
|\alpha\rangle_{[n+1:N]},
\label{eq:ch2-schmidt}
\end{equation}
其中 $r_n=\mathrm{rank}(M^{(n)})$，$\lambda_{\alpha}^{(n)}\ge 0$，且 Schmidt 系数按从大到小排序。相应的二分纠缠熵为
\begin{equation}
S_n=-\sum_{\alpha}(\lambda_{\alpha}^{(n)})^2\log (\lambda_{\alpha}^{(n)})^2.
\label{eq:ch2-entanglement-entropy}
\end{equation}
若奇异值谱在较小的 $\alpha$ 处已经明显衰减，只保留前 $\chi_n$ 个主分量就能得到一条低秩近似，其截断误差为
\begin{equation}
\varepsilon_n^2=\sum_{\alpha>\chi_n}(\lambda_{\alpha}^{(n)})^2.
\label{eq:ch2-truncerr}
\end{equation}
式 \eqref{eq:ch2-truncerr} 给出“丢弃权重”的定义。对应的数值操作是保留当前动力学中贡献最大的纠缠通道，并把权重很小的 Schmidt 分量从数值表示中截去。

上述结构也解释了为什么 MPS 在一维问题上特别有效。对一维有能隙体系的基态，面积律表明纠缠熵不会随着系统长度线性增长，因此存在高效的 MPS 逼近 \cite{Hastings2007AreaLaw}。所研究链路属于非平衡动力学问题，严格的面积律不能直接搬用；但发射、QFC 和记忆核都属于有限时间内的顺序相互作用，纠缠增长主要沿着因果锥逐步传播，因此在有限时间窗和有限带宽条件下，真正需要保留的键维通常仍远小于形式上的指数上界 \cite{Vidal2003SlightlyEntangled,Vidal2004TEBD,Pichler2016TimeDelaysMPS}。时间 bin 方法因此能够在量子光学延迟链路中保持可计算性。

\subsection{规范形式和可观测量收缩}[Canonical Forms and Observable Contraction]

把式 \eqref{eq:ch2-full-state} 逐站点做低秩分解后，可得到矩阵乘积态表示
\begin{equation}
C_{s_1\cdots s_N}
=
\sum_{\alpha_1,\ldots,\alpha_{N-1}}
A^{[1]s_1}_{\alpha_1}
A^{[2]s_2}_{\alpha_1\alpha_2}
\cdots
A^{[N]s_N}_{\alpha_{N-1}},
\label{eq:ch2-mps}
\end{equation}
其中 $\alpha_k=1,\ldots,\chi_k$ 是虚拟指标，$\chi=\max_k\chi_k$ 为键维上界。MPS 的参数量由 $d^N$ 降为 $O(Nd\chi^2)$，指数复杂度被压缩成对键维的多项式复杂度 \cite{Schollwoeck2011MPS,Orus2014TensorNetworks}。

MPS 最重要的一个直觉解释，是把虚拟指标看成“相关信息在链上的输运通道”。若某条键上的 Schmidt 谱很窄，那么该处真正需要承载的信息就不多；若某条键上的 Schmidt 谱突然变宽，通常说明相互作用在该处生成了更强的关联。于是键维成为当前动力学复杂度的一种直接刻画。这个解释对理解后续参数扫描很有帮助：当窗口加宽、滤波记忆增强或 bin 数增多时，计算成本上升往往首先反映为某几条键上的谱变宽，链上其他位置未必同步变复杂。

同一个物理态对应无穷多组 MPS 张量，因此还需要利用规范自由度把表示整理成数值上更稳定的形式。相邻张量之间的规范变换写成
\begin{align}
A^{[n]} &\rightarrow A^{[n]}X,\\
A^{[n+1]} &\rightarrow X^{-1}A^{[n+1]},
\label{eq:ch2-gauge-transform}
\end{align}
其中 $X$ 为任意可逆矩阵。利用这一定义，可以把 MPS 写成左规范、右规范或混合规范形式。混合规范最常用，因为某一条中心键上的对角矩阵就直接给出该切分处的 Schmidt 系数。采用混合规范后有两点直接好处：第一，局域门更新后的重分解只影响附近少数张量；第二，期望值、关联函数和概率收缩都可以借助左右环境递推稳定计算 \cite{Schollwoeck2011MPS,Orus2014TensorNetworks}。

设 $|\psi\rangle$ 已处于适当规范形式，则单体算符 $\hat O_n$ 的期望值写为
\begin{equation}
\langle \hat O_n \rangle
=
\frac{\langle \psi|\hat O_n|\psi\rangle}{\langle \psi|\psi\rangle},
\label{eq:ch2-onebody-expectation}
\end{equation}
它可通过左环境与右环境的逐步收缩得到。两点关联函数的计算完全类似，只是在对应位置插入两个算符核。对端到端量子接口问题而言，后续真正需要的量并不局限于普通期望值，还包括成功事件概率、Bell 态保真度和条件后的原子态密度矩阵。它们的共同点，是最终都能归约为“在 MPS 链上插入有限个局域算符，再做一次可控收缩”。因此，MPS 在这里不只是压缩存储工具，也是整个概率计算接口的底层结构。

\begin{figure}[htbp]
  \centering
  \includegraphics[width=0.88\textwidth]{fig_mps_chain.pdf}
  \caption{所研究链路在 MPS 中的站点排布示意。发射体、时间 bin 和滤波记忆模都被编码为一维链上的局域站点；活跃作用区域通过门调度在链上逐步移动。该排布把连续时间量子光学问题转化为顺序局域更新问题。}
  \label{fig:ch2-mpschain}
\end{figure}

若完全在态端维护密度矩阵，局域维数会进一步平方，计算成本显著上升。对开放系统而言，矩阵乘积密度算符（matrix product density operator, MPDO）和超算符 TEBD 是一条标准路线 \cite{Verstraete2004MPDO,Zwolak2004MixedStateDynamics}。所研究链路把必须显式保留的相干历史尽量留在纯态 MPS 中，把线性光学和探测噪声统计放到 effect 端处理。这一选择首先有利于控制算力开销，同时也贴合链路本身的物理结构：主要相干历史集中在发射与有限带宽记忆之中，探测端更适合表述为条件概率与广义测量问题。

\section{TEBD 时间推进}[Time-Evolving Block Decimation]

MPS 解决的是“如何表示态”，TEBD 解决的是“如何更新态” \cite{Vidal2004TEBD,Paeckel2019TimeEvolutionMPS}。在所研究链路里，时间推进通过发射门、QFC 门、滤波记忆门和必要的 SWAP 门组成的局域门时间表来实现，并在每一步后恢复适合收缩的规范形式。每一个数值操作都对应一个明确的物理步骤，因果顺序和误差来源也因此更容易检查 \cite{Vidal2004TEBD,Paeckel2019TimeEvolutionMPS}。

\subsection{局域门更新}[Local Gate Updates]

单站点门 $U^{[n]}$ 只作用于一个物理腿，其更新规则为
\begin{equation}
\tilde A^{[n]s'}_{\alpha\beta}
=
\sum_s U^{[n]}_{s's}\,A^{[n]s}_{\alpha\beta}.
\label{eq:ch2-one-site-update}
\end{equation}
双站点门则更为关键，因为绝大多数纠缠增长都来自两站点相互作用。把相邻站点 $n,n+1$ 合并为二站点张量 $\Theta$ 后，可写成
\begin{equation}
\Theta
\xrightarrow{\,U^{[n,n+1]}\,}
\tilde\Theta
\xrightarrow{\mathrm{SVD}}
\tilde A^{[n]}\tilde\Lambda^{[n]}\tilde A^{[n+1]},
\label{eq:ch2-two-site-update}
\end{equation}
然后按式 \eqref{eq:ch2-truncerr} 对过小奇异值进行截断。这个步骤之所以重要，是因为“施加局域门”和“恢复低秩表示”是两件不同的事。前者反映物理相互作用，后者负责把更新后的态重新压回可计算流形。若没有第二步，键维会随着时间步迅速失控；若第二步做得过于激进，又会把小概率通道系统性低估掉。

对所研究链路而言，式 \eqref{eq:ch2-two-site-update} 对应每一类局域物理过程的统一模板。发射阶段使用“发射体 $+$ 当前时间 bin”两站点门；QFC 阶段使用“当前 bin 内不同频率模式”形成的单站点或嵌入式局域门；滤波阶段使用“当前 bin $+$ 记忆模”两站点门；为了让后续 bin 与记忆模相邻，还需要插入 SWAP 门。不同物理过程虽然矩阵元素不同，数值骨架却保持一致，整个算法始终围绕局域门更新、局域 SVD 和局域重规范三件事情展开。

\subsection{时间推进调度}[Time-Evolution Schedule]

当系统哈密顿量可写为最近邻和 $H=\sum_n h_{n,n+1}$ 时，TEBD 通常利用奇偶键分组 $H=H_{\mathrm{o}}+H_{\mathrm{e}}$，二阶 Suzuki--Trotter 近似为
\begin{equation}
e^{-iH\Delta t}
=
e^{-iH_{\mathrm{o}}\Delta t/2}
e^{-iH_{\mathrm{e}}\Delta t}
e^{-iH_{\mathrm{o}}\Delta t/2}
+O(\Delta t^3).
\label{eq:ch2-trotter2}
\end{equation}
从总演化时间 $T$ 来看，二阶 Trotter 误差通常按 $O(T\Delta t^2)$ 累积 \cite{Vidal2004TEBD,Paeckel2019TimeEvolutionMPS}。步长足够小时，Trotter 误差会减小；不过实际总误差还会叠加每一步的 Schmidt 截断和局域维度截断。因此，步长 $\Delta t$、最大键维 $\chi_{\max}$ 和局域空间截断需要联合做收敛检查，不能分别凭经验设定。

在所研究链路里，TEBD 的实际调度比式 \eqref{eq:ch2-trotter2} 稍微复杂一些，但思想完全相同：
\begin{enumerate}
  \item 在发射阶段，发射体按时间顺序依次与各个 bin 发生局域相互作用；
  \item 在 QFC 阶段，对每个 bin 施加频率转换门，并把转换后的频谱信息保留在 bin 的局域基中；
  \item 在滤波阶段，让当前 bin 与记忆模相互作用，再用 SWAP 门把记忆模移动到下一个待作用位置附近；
  \item 在测量阶段，计算流程从态端门更新切换到 effect 收缩。
\end{enumerate}
这种调度方式看上去比“统一哈密顿量时间推进”更工程化，但它更接近真实实验链路的操作顺序，也更容易和代码中的门矩阵一一对应起来。

\begin{figure}[htbp]
  \centering
  \includegraphics[width=0.90\textwidth]{fig_tebd_pattern.jpg}
  \caption{局域门调度示意。一个时间步内的主要工作包括相邻站点门作用、奇异值分解回写以及必要的 SWAP 调度。对于所研究链路，发射门、QFC 门和滤波记忆门都被写成这一调度模式中的局域更新。}
  \label{fig:ch2-tebd}
\end{figure}

从复杂度角度看，在局域维数近似均匀时，双站点更新的主开销来自局域 SVD，单步成本可估计为
\begin{equation}
O(Nd^3\chi^3),
\label{eq:ch2-single-step-cost}
\end{equation}
总成本约为 $O(TNd^3\chi^3/\Delta t)$。这个估计虽然粗糙，却已经足够解释后续计算中的主要瓶颈：bin 数增加会线性拉长链长 $N$，更细的时间分辨率会增加 $T/\Delta t$，而滤波记忆和更强的不完美会抬高所需键维 $\chi$。其中最敏感的通常仍是 $\chi$，因为它以三次方进入成本估计。后续收敛性检查通常先看哪些位置的 Schmidt 谱确实要求更高的键维，再决定是否需要继续增加时间 bin。

在时间演化方法的选择上，TEBD 只是可选方案之一。Paeckel 等人的综述比较了 TEBD、Krylov 方法和时间依赖变分原理（TDVP）等多种 MPS 时间推进策略 \cite{Paeckel2019TimeEvolutionMPS}。当前链路最终选择 TEBD，主要因为局域门来源清晰、顺序相互作用强、且每一步都能和实验物理过程建立直接对应。对以“门序列”组织的接口问题而言，这种可解释性往往比抽象的全局推进更重要。

\section{条件测量的算符方法}[Operator Method for Conditional Measurements]

量子接口链路末端给出的结果是一份带条件的实验记录：某个时间窗里哪些探测器点击了、点击模式是否属于成功集合、成功后两个远端原子处于什么条件态。只用态矢量演化并不能直接回答这些问题，因此需要把开放系统和广义测量接到同一条概率链上。Kraus 表示给出开放系统信道的离散表述，POVM 给出广义测量的记录空间，Heisenberg 对偶映射则把二者连接起来，使线性光学和探测端能够在 effect 端统一处理 \cite{Kraus1983StatesEffectsOperations,Lindblad1976Semigroups,GoriniKossakowskiSudarshan1976,NielsenChuang2010QCQI}。

\subsection{Kraus 表示}[Kraus Representation]

Markov 近似下的开放系统动力学通常由 GKSL 主方程描述：
\begin{equation}
\dot\rho=-i[H,\rho]+\sum_{\mu}\left(L_{\mu}\rho L_{\mu}^{\dagger}-\frac{1}{2}\{L_{\mu}^{\dagger}L_{\mu},\rho\}\right),
\label{eq:ch2-gksl}
\end{equation}
其中 $H$ 负责相干演化，$L_{\mu}$ 是各个耗散通道的跃迁算符。对一个足够短的时间步长 $\Delta t$，式 \eqref{eq:ch2-gksl} 可以写成完全正且保持迹映射
\begin{equation}
\mathcal E(\rho)=\sum_{\mu}K_{\mu}\rho K_{\mu}^{\dagger},
\label{eq:ch2-kraus}
\end{equation}
并满足
\begin{equation}
\sum_{\mu}K_{\mu}^{\dagger}K_{\mu}=I.
\label{eq:ch2-kraus-complete}
\end{equation}
Kraus 表示的价值在于把“开放系统演化”写成一族可与测量记录对应的分支过程。对短时间步，常用的一组近似 Kraus 算符可写为
\begin{equation}
K_0=I-iH_{\mathrm{eff}}\Delta t+O(\Delta t^2),
\qquad
K_{\mu}=\sqrt{\Delta t}\,L_{\mu},
\label{eq:ch2-short-time-kraus}
\end{equation}
其中 $H_{\mathrm{eff}}=H-\frac{i}{2}\sum_{\mu}L_{\mu}^{\dagger}L_{\mu}$ 是非厄米有效哈密顿量。于是“无跳跃”与“发生某一跃迁”便自然对应到不同的分支。

若把单次演化看成一条随机轨迹，则纯态在第 $\mu$ 个分支上的更新为
\begin{equation}
|\psi\rangle\rightarrow |\psi_{\mu}\rangle=
\frac{K_{\mu}|\psi\rangle}{\|K_{\mu}|\psi\rangle\|},
\qquad
p_{\mu}=\|K_{\mu}|\psi\rangle\|^2.
\label{eq:ch2-trajectory}
\end{equation}
对所有轨迹做统计平均后，又会回到式 \eqref{eq:ch2-kraus} 的密度矩阵描述。这种“单条轨迹是纯态、平均后是开放系统”的双重表述，在量子光学里尤其自然，因为实验记录本来就以点击事件或跃迁事件的形式出现 \cite{DalibardCastinMolmer1992QuantumJumps,Daley2014QuantumTrajectories}。对于所研究链路，发射体与时间 bin 的相互作用、自由空间泄漏与部分器件损耗都可以在这一框架中得到统一描述。

\subsection{POVM 对偶映射}[POVM and Adjoint Mapping]

广义测量由一组正算符 $\{\Pi_m\}$ 描述，它们满足
\begin{align}
\Pi_m &\succeq 0,\\
\sum_m \Pi_m &= I,
\label{eq:ch2-povm-complete}
\end{align}
并给出记录 $m$ 的概率
\begin{equation}
p_m=\mathrm{Tr}(\Pi_m\rho).
\label{eq:ch2-povm-prob}
\end{equation}
若写成测量算符形式 $\Pi_m=M_m^{\dagger}M_m$，则条件后的后验态为
\begin{equation}
\rho_m=\frac{M_m\rho M_m^{\dagger}}{\mathrm{Tr}(M_m\rho M_m^{\dagger})}.
\label{eq:ch2-posterior-single}
\end{equation}
对当前问题而言，单个记录 $m$ 可以理解为一个探测时间窗内的四路点击模式及其标签；真正关心的成功事件通常对应一组记录的并集，例如可分辨的 partial BSM 成功模式集合。

关键一步在于对偶映射。对任意信道 $\Lambda$ 及任意 effect $E$，有
\begin{equation}
\mathrm{Tr}[\Lambda(\rho)E]
=
\mathrm{Tr}[\rho\,\Lambda^{\dagger}(E)],
\label{eq:ch2-adjoint}
\end{equation}
其中
\begin{equation}
\Lambda^{\dagger}(E)=\sum_k K_k^{\dagger}EK_k.
\label{eq:ch2-adjoint-kraus}
\end{equation}
若 $\Lambda$ 是幺正门 $U$，则 $\Lambda^{\dagger}(E)=U^{\dagger}EU$；若 $\Lambda$ 是损耗、暗计数或线性光学网络对应的信道，则 effect 会在输入端变成一组新的等效算符。于是，探测器可以不在态端显式模拟，概率计算时直接作为输入端的等效测量算符出现 \cite{Kraus1983StatesEffectsOperations,NielsenChuang2010QCQI}。

将其写为全链路形式。设从输入端到探测端依次经过信道 $\Phi_1,\Phi_2,\ldots,\Phi_n$，探测端记录 $m$ 的 effect 为 $F_m$，则有
\begin{equation}
\mathrm{Tr}\!\left[F_m(\Phi_n\circ\cdots\circ\Phi_1)(\rho_{\mathrm{in}})\right]
=
\mathrm{Tr}\!\left[\bigl((\Phi_1^{\dagger}\circ\cdots\circ\Phi_n^{\dagger})(F_m)\bigr)\rho_{\mathrm{in}}\right].
\label{eq:ch2-op-impl-consistency}
\end{equation}
式 \eqref{eq:ch2-op-impl-consistency} 说明，只要链路中的这些器件不反过来影响上游的发射动力学，那么就可以把它们全部压缩为输入端的等效 effect，而无需在态端显式传播整个探测网络。在该链路中，这正是“发射与滤波记忆留在 state side，光纤、分束器和探测器放到 effect side”的理论依据。

\begin{figure}[htbp]
  \centering
  \includegraphics{fig_tikz_ch2_duallane}
  \caption{所研究链路的双通道方法示意。状态侧保留会积累相干历史的发射和有限带宽记忆，测量侧则把光纤、线性光学和探测器统一写成等效 effect 并向输入端推回。该分工在保持物理解释性的同时避免混态维数在态端全面膨胀。}
  \label{fig:ch2-duallane}
\end{figure}

对于成功事件集合 $\mathcal S$，条件概率和条件态可写为
\begin{equation}
p_{\mathcal S}=\sum_{m\in\mathcal S}p_m,
\label{eq:ch2-success-prob}
\end{equation}
\begin{equation}
\rho_{\mathcal S}=
\frac{1}{p_{\mathcal S}}\sum_{m\in\mathcal S}M_m\rho M_m^{\dagger}.
\label{eq:ch2-success-state}
\end{equation}
Bell 态保真度、成功率、真假成功分解以及 CHSH 指标都建立在式 \eqref{eq:ch2-success-prob} 和式 \eqref{eq:ch2-success-state} 之上。所有“事件口径”的差异都能追溯到 $\mathcal S$ 或 $F_m$ 的定义，态演化本身与事件定义也能保持分离。

\section{观测量定义及误差控制}[Observables and Error Control]

在端到端接口仿真中，观测量定义和误差判据需要与演化方程同时给出。在完成 MPS/TEBD 与算符框架之后，还需要把链路中核心观测量和收敛检查方法明确下来：一是固定概率与条件态的符号和口径；二是区分“可计算”与“结果可信”对应的误差边界。

\subsection{主要观测量}[Primary Observables]

在中间站 Bell 测量链路中，最先需要区分的是“到达了什么”和“宣告了什么”。前者由输入端光子占据情况刻画，例如 $p_{11}$ 表示两臂各有一个光子到达中间站的概率；后者由成功记录集合 $\mathcal S$ 决定，例如 $p_s$ 表示所有成功 click pattern 的总概率。若再把 success effect 分解为真成功与假成功两部分，就可进一步定义
\begin{equation}
p_s = \sum_{m\in\mathcal S}\mathrm{Tr}(E_m\rho),
\qquad
p_t = \sum_{m\in\mathcal S}\mathrm{Tr}(E_m^{\mathrm{true}}\rho),
\qquad
p_f = p_s-p_t,
\label{eq:ch2-true-false}
\end{equation}
以及在双臂单光子真正到达条件下的真成功率
\begin{equation}
p_{t|11}=\frac{p_t}{p_{11}}.
\label{eq:ch2-conditional-true}
\end{equation}
这里 $E_m^{\mathrm{true}}$ 表示那些确实由目标物理过程引发的成功 effect，$E_m^{\mathrm{false}}=E_m-E_m^{\mathrm{true}}$ 则吸收由多光子、背景、暗计数等造成的假宣告成分。这一分解对实验解释很重要，因为总成功率变高并不必然意味着接口质量变好，增加的也可能主要是假成功通道。

条件态质量通常用目标 Bell 态保真度来衡量。若目标态记为 $|\mathcal B\rangle$，则有
\begin{equation}
F_{\mathcal B}=\langle \mathcal B|\rho_{\mathcal S}|\mathcal B\rangle.
\label{eq:ch2-bell-fidelity}
\end{equation}
当 $\rho_{\mathcal S}$ 已被压缩到两原子子空间后，还可以引入相关张量
\begin{equation}
T_{ij}=\mathrm{Tr}\left[\rho_{\mathcal S}(\sigma_i\otimes\sigma_j)\right],
\qquad i,j\in\{x,y,z\},
\label{eq:ch2-correlation-tensor}
\end{equation}
并用 Horodecki 判据给出 CHSH 最大值
\begin{equation}
S_{\max}=2\sqrt{\lambda_1+\lambda_2},
\label{eq:ch2-chsh-max}
\end{equation}
其中 $\lambda_1,\lambda_2$ 为矩阵 $T^{\mathsf T}T$ 的最大两个特征值 \cite{Horodecki1995CHSH}。于是，保真度回答“条件态离目标 Bell 态有多近”，$S_{\max}$ 回答“该条件态是否仍能支持非定域关联”。二者并不冗余：同一保真度下，不同的噪声结构可能给出不同的 CHSH 表现。

两光子不可区分性则常用 Hong--Ou--Mandel 可见度表征。若平行偏振与正交偏振配置下的符合计数分别记为 $C_{\parallel}$ 与 $C_{\perp}$，则
\begin{equation}
V_{\mathrm{HOM}}=\frac{C_{\perp}-C_{\parallel}}{C_{\perp}}.
\label{eq:ch2-hom-visibility}
\end{equation}
该量并不直接等于最终远程原子纠缠保真度，但它对时间抖动、频率失配和多模污染高度敏感，因此在接口标定中具有前置诊断价值 \cite{HongOuMandel1987,vanleent2022entangling}。

\begin{table}[htbp]
  \centering
  \caption{主要观测量及定义}
  \label{tab:ch2-observables}
  \begin{tabular}{p{2.2cm}p{3.0cm}p{8.8cm}}
    \toprule
    符号 & 名称 & 定义或物理含义 \\
    \midrule
    $p_{11}$ & 双臂单光子到达概率 & 两臂各有一个光子到达中间站输入端的概率，是判断 BSM 是否具备理想输入条件的基线量。 \\
    $p_s$ & 宣告成功概率 & 所有成功 click pattern 概率之和，反映实验上“看到一次成功事件”的总频率。 \\
    $p_t$ & 真成功概率 & 成功事件中真正由目标单光子干涉过程产生的部分。 \\
    $p_f$ & 假成功概率 & 由多光子、背景、暗计数等引发的成功宣告成分，满足 $p_f=p_s-p_t$。 \\
    $p_{t|11}$ & 条件真成功率 & 在理想输入条件 $11$ 到达下的真成功率，用于把输入制备质量和 BSM 判别质量分开。 \\
    $F_{\mathcal B}$ & Bell 态保真度 & 条件后原子态对目标 Bell 态 $|\mathcal B\rangle$ 的保真度。 \\
    $S_{\max}$ & CHSH 最大值 & 条件态在最优测量基下的 CHSH 指标，用于衡量非定域关联。 \\
    $V_{\mathrm{HOM}}$ & HOM 可见度 & 两光子不可区分性的前置诊断量，对时间失配和谱失配尤为敏感。 \\
    \bottomrule
  \end{tabular}
  \begin{minipage}{0.94\textwidth}
    \footnotesize 注：表中量的概率定义基于 Kraus 表示与 POVM 条件化整理，参见 \cite{Kraus1983StatesEffectsOperations,NielsenChuang2010QCQI}；$S_{\max}$ 的两比特判据见 \cite{Horodecki1995CHSH}；HOM 可见度的标准定义可追溯至 \cite{HongOuMandel1987}。
  \end{minipage}
\end{table}

\subsection{误差来源}[Error Sources]

端到端接口仿真的误差来源并不只有一类。最容易想到的是 Trotter 误差和 Schmidt 截断误差，但对当前问题来说，至少还应额外看清楚三件事：连续时间离散化本身带来的时间分箱误差；局域基空间有限维近似带来的模型截断；以及把光纤和探测网络推到 effect 端之后所对应的适用边界。若这些边界没有被明确写出，数值结果即便收敛，也很难说清“究竟收敛到了哪个模型”。

首先是时间分箱误差。离散化之后，动力学中的关键无量纲组合通常以 $g\sqrt{\Delta t}$、$\kappa\Delta t$、$\gamma\Delta t$ 这类形式出现。若 $\Delta t$ 过大，短时近似会失效，波包形状、到达概率和 HOM 曲线都会出现明显的网格依赖。判断这一误差时，需要同时比较波包包络、成功率和保真度在 $\Delta t\to \Delta t/2$ 时是否稳定，而不能只盯住某一个单独概率。只有当多种指标同时不再明显漂移时，离散时间步长才算真正收敛。

其次是 Schmidt 截断与局域维度截断。前者反映“保留多少纠缠通道”，后者反映“每个站点允许出现哪些局域物理状态”。在所研究实现中，时间 bin 被截断到真空加单光子子空间，滤波记忆模也只保留有限维数。这一做法的前提，是多光子占据和更高阶记忆激发在目标参数区间内确实可以忽略。若背景噪声、泵浦泄漏或强驱动使这些高阶占据不再小，那么即便 $\chi$ 已经收敛，模型本身仍可能需要扩充局域基。换言之，“键维收敛”和“局域空间足够大”是两类不同问题，不能相互替代 \cite{Schollwoeck2011MPS,Paeckel2019TimeEvolutionMPS}。

再次是 effect 端推回的模型边界。把分束器、光纤损耗和探测器响应统一压缩成输入端 effect 的前提，是这些过程不会反向影响上游的发射和记忆动力学，也不会在传播途中引入必须显式保存在态端的额外相干历史。对线性损耗、静态相位、探测效率和暗计数而言，这一处理是自然的；但若传播过程本身产生显著的多模混合、反馈或与上游强相关的记忆效应，就需要把更多传播物理显式纳入 state side，而不能全部交给 effect 端吸收 \cite{Pichler2016TimeDelaysMPS,Wang2017MPSCascadedQuantumSystems,ArranzRegidor2021WaveguideQED}。

\begin{table}[htbp]
  \centering
  \caption{端到端量子接口仿真的主要误差来源和控制参数}
  \label{tab:ch2-errors}
  \begin{tabular}{p{2.6cm}p{2.6cm}p{4.0cm}p{4.8cm}}
    \toprule
    误差来源 & 主要控制参数 & 诊断量 & 物理含义 \\
    \midrule
    时间分箱误差 & $\Delta t$、总 bin 数 $N_t$ & 波包形状、$p_s$、$F_{\mathcal B}$ 对步长的收敛性 & 连续时间场被离散为有限时间窗后产生的近似。 \\
    Trotter 误差 & $\Delta t$、Trotter 阶数 & 同一参数下不同步长结果是否一致 & 局域门乘积替代真实时间序演化带来的系统误差。 \\
    Schmidt 截断误差 & $\chi_{\max}$、丢弃权重 $\varepsilon_n^2$ & 键谱尾部、核心观测量对 $\chi_{\max}$ 的稳定性 & 为控制计算规模而丢弃较小 Schmidt 分量的误差。 \\
    局域空间截断 & bin 维数、记忆模维数 & 高阶占据是否可忽略、概率是否出现异常泄漏 & 有限维站点是否足以容纳实际局域物理过程。 \\
    effect 端模型边界 & 推回器件的范围、传播模型假设 & 与更显式模型或简化基准的比较 & 判断哪些传播和探测过程可以安全地压缩为等效 effect。 \\
    统计误差 & 枚举范围、抽样次数 & 重复运行波动、置信区间 & 由记录空间抽样或有限样本统计带来的不确定度。 \\
    \bottomrule
  \end{tabular}
  \begin{minipage}{0.94\textwidth}
    \footnotesize 注：表中误差分类综合了 MPS/TEBD 时间推进、开放系统量子轨迹以及时间延迟光子线路的常见误差口径，参见 \cite{Vidal2004TEBD,Schollwoeck2011MPS,Paeckel2019TimeEvolutionMPS,Daley2014QuantumTrajectories,Pichler2016TimeDelaysMPS,Wang2017MPSCascadedQuantumSystems}。
  \end{minipage}
\end{table}

对所研究链路而言，更实用的收敛顺序通常如下：先用波包与 HOM 诊断确定 $\Delta t$ 已足够小，再检查核心观测量对 $\chi_{\max}$ 的依赖，随后确认局域维度截断与窗口设定不会改变成功口径，最后再讨论更细的探测统计与抽样误差。原因很直接。若离散时间步长本身还没有稳定下来，继续提高键维通常只是更精确地逼近一个尚未收敛的离散模型；若时间网格已经可靠，而关键观测量仍然随 $\chi_{\max}$ 漂移，问题多半出在纠缠压缩，而非物理参数本身。

\section{小结}[Chapter Summary]

端到端量子接口链路的方法框架可概括为：连续时间光场离散为时间 bin 后形成一维链式排布；利用 Schmidt 分解与 MPS 规范形式将指数复杂度压缩到由键维控制的低秩流形；利用 TEBD 将发射、QFC、滤波记忆和 SWAP 调度统一为局域门序列；借助 Kraus 表示、POVM 与对偶映射把开放系统噪声和条件测量组织到同一条概率链中。该框架中，state side 显式保留发射与有限带宽记忆，effect side 负责线性光学与探测统计；成功率、真假成功分解、Bell 态保真度、CHSH 指标与 HOM 可见度在统一算符框架下定义，时间分箱误差、Trotter 误差、Schmidt 截断和 effect 端模型边界共同构成收敛检查。




